%--------------%
%->> Section 4.2: 系统实现
%------------%

\section{系统实现}

本节进一步说明第三章提出的三类方法在平台中的具体落地方式。围绕故障指纹构建、SOP自动化构建和在线故障诊断三类核心方法，平台分别形成了对应的实现链路。表\ref{tab:chap3_chap4_impl}给出了第三章方法与平台实现之间的对应关系。

\begin{table}[htbp]
    \centering
    \bicaption{\enspace 第三章方法与平台实现的对应关系}{\enspace Mapping of Chapter 3 Methods to Platform Components}
    \label{tab:chap3_chap4_impl}
    \begingroup
    \small
    \setlength{\tabcolsep}{4pt}
    \renewcommand{\arraystretch}{1.25}
    \begin{tabularx}{\textwidth}{>{\raggedright\arraybackslash}p{3.3cm}>{\raggedright\arraybackslash}p{4.0cm}>{\raggedright\arraybackslash}X}
        \toprule
        \textbf{第三章方法} & \textbf{平台中的主要位置} & \textbf{具体实现} \\
        \midrule
        故障指纹构建方法 & 故障指纹库、指纹管理模块、编码模型服务 & 负责故障指纹编码、存储、聚类维护和在线检索，并返回相似故障及其绑定的SOP \\
        SOP 自动化构建方法 & 工作区、\texttt{Exploration-Skill}、\texttt{Rule-Skill}、轨迹存储 & 负责生成探索轨迹、提取决策单元，并完成SOP更新 \\
        在线故障诊断方法 & \texttt{Diagnosis-Skill}、观察 subagent、结果回流更新机制 & 负责读取告警与SOP、收集证据、推进诊断，并将结果回流至后续更新流程 \\
        \bottomrule
    \end{tabularx}
    \endgroup
\end{table}

\subsection{故障指纹构建方法的实现}

第 3.2 节提出的传播感知多模态故障指纹构建方法，主要由故障指纹库、指纹管理模块和编码模型服务共同实现。由于该部分任务以批处理和持续维护为主，涉及特征编码、向量存储与聚类更新，因此其核心计算逻辑部署在后端服务侧。

在离线阶段，系统首先从历史故障实例中抽取指标、日志和链路等多模态数据，再调用编码模型服务生成 MMF、CTF、LCF、LSF 和 CMF 五类特征，并将其拼接形成故障指纹，存入故障指纹库。除保存指纹向量外，故障指纹库还维护聚类中心以及故障簇与SOP的绑定关系。指纹管理模块负责这些记录的更新，并定期对已有指纹执行DBSCAN聚类，以刷新聚类中心和绑定结果。

在在线阶段，系统接收到告警后，由后端首先调用指纹管理模块完成相似故障检索，返回最近邻故障簇及其绑定的SOP，再将检索结果同步到当前工作区，作为 \texttt{Diagnosis-Skill} 的初始上下文。通过这一实现方式，在线诊断在开始时即可获得与当前告警最接近的历史故障模式及其对应的排查流程。

\subsection{SOP自动化构建方法的实现}

第 3.3 节提出的SOP自动化构建方法，在平台中由 \texttt{Exploration-Skill} 和 \texttt{Rule-Skill} 依次完成。前者负责围绕给定故障场景生成可分析的探索轨迹，后者负责从这些轨迹中提取决策单元并更新 SOP 知识库。之所以将二者分开执行，是因为探索阶段关注证据扩展和路径尝试，而规则阶段关注结构整理与知识沉淀，两类任务在输入、输出和控制目标上均存在明显差异。

在具体实现中，后端首先围绕某一故障实例创建工作区，并初始化故障类型先验、相关服务范围和当前SOP。随后，系统启动多个代理，基于 \texttt{Exploration-Skill} 并行执行故障探索任务，从而生成多条 SOP 执行轨迹。在每个代理的执行过程中，\texttt{Exploration-Skill} 负责约束探索行为，包括下一步检查对象、路径推进方式以及当前探索的结束条件；若某一步需要补充指标、日志或链路证据，则创建观察 subagent，由其通过 \texttt{Observe MCP} 获取所需信息并返回结果。

探索结束后，生成的轨迹被保存至轨迹存储。随后系统启动 \texttt{Rule-Skill}，读取这些轨迹，对其中的关键判断、跳转条件和结论形成过程进行拆解，完成决策单元提取、语义对齐和权重计算。最后，\texttt{Rule-Skill} 根据高权重决策单元对原有 SOP 执行节点的增删改操作，并将更新后的有向图同步至SOP知识库。由此，SOP的更新过程被组织为“探索生成轨迹—规则整理轨迹—知识库完成更新”的连续实现链路。

\subsection{在线故障诊断方法的实现}

第 3.4 节提出的SOP驱动多智能体在线故障诊断方法，主要由 \texttt{Diagnosis-Skill}、观察 subagent 和结果回流更新机制共同实现。系统接收到告警后，后端首先完成两项准备工作：一是调用故障指纹检索接口，获得最近邻故障簇及其绑定的SOP；二是创建工作区，并将告警信息、检索结果和选定的SOP初始化到当前任务上下文中。随后，OpenCode 启动诊断会话，加载 \texttt{Diagnosis-Skill}，并据此初始化焦点与感知域。

诊断开始后，\texttt{Diagnosis-Skill} 负责主流程推进。它根据当前焦点和已有证据决定下一步动作，并在运行过程中执行三类核心操作：\texttt{navigate} 表示沿当前 SOP 继续推进，\texttt{escape} 表示在现有SOP无法覆盖当前场景时通过逃逸机制扩展新的诊断路径，\texttt{conclude} 表示输出当前诊断结论。上述操作由 \texttt{Diagnosis MCP} 提供相应的接口和执行约束。

与离线探索阶段相同，\texttt{Diagnosis-Skill} 并不直接访问指标、日志和链路数据。当诊断过程需要补充某一焦点范围内的可观测证据时，主会话会创建观察 subagent，由后者通过 \texttt{Observe MCP} 查询所需数据，再将结果返回主流程。这样，诊断过程中的路径控制与证据采集被明确分离：前者由主会话负责，后者由观察子代理承担。

诊断结束后，系统将诊断结论、关键证据和完整执行过程保存为候选诊断记录。经过人工审核后，相关结果将进一步进入故障指纹更新流程和SOP增量更新流程，用于补充后续任务所依赖的历史知识。通过这一实现方式，在线诊断不仅能够利用已有知识完成当前故障定位，也能够将新的有效经验持续回流至平台知识体系中。

综上，第三章提出的三类方法在平台中分别落地为故障指纹检索、SOP自动化构建和在线诊断执行三条实现链路。三者侧重点不同，但通过工作区、知识库以及结果回流更新机制连接在一起，共同构成了平台的核心实现流程。